% Appendix A — Notation and Conventions
\chapter{Notation and Conventions}
\label{app:notation}

This appendix collects the notational conventions and symbol definitions used
throughout the book, providing a single reference point for readers navigating
between chapters.

\section{Index Conventions}
\label{sec:not-indices}

Greek indices $\mu, \nu, \ldots$ run over the four spacetime dimensions
$0, 1, 2, 3$. Latin indices $i, j, \ldots$ run over the three spatial
dimensions $1, 2, 3$. Repeated indices imply Einstein summation unless stated
otherwise.

\section{Metric Signature}
\label{sec:not-signature}

We adopt the $({-}{+}{+}{+})$ signature throughout, so that time-like vectors
have negative norm $g_{\mu\nu} V^\mu V^\nu < 0$.

\section{Natural Units}
\label{sec:not-units}

We work in geometrised units with $G = c = 1$, so that mass, length, and time
all have dimensions of length. Factors of $G$ and $c$ are restored explicitly
where needed for dimensional clarity or astrophysical estimates.

\section{Symbol Table}
\label{sec:not-symbols}

% TODO: Populate with a comprehensive tabular of symbols as chapters are
% written. Each entry: symbol, definition, first-introduced chapter.
A comprehensive symbol table will be provided here, listing each mathematical
symbol alongside its definition and the chapter where it is first introduced.

\paragraph{Key References.}
Our conventions follow \cite{MTW1973} for index notation and
\cite{Wald1984} for the metric signature; geometrised units are standard in
\cite{Baumgarte2010}.
